% Generated draft: requires mathematical and source review.
\documentclass[11pt]{article}
\usepackage[margin=25mm]{geometry}
\usepackage{fontspec}
\setmainfont{Latin Modern Roman}
\usepackage{amsmath,amssymb,mathtools,unicode-math}
\setmathfont{Latin Modern Math}
\usepackage{enumitem,needspace,xurl,hyperref}
\setlength{\parindent}{0pt}
\setlength{\parskip}{5pt}
\setlength{\emergencystretch}{3em}
\newcommand{\N}{\mathbb{N}}
\newcommand{\Z}{\mathbb{Z}}
\newcommand{\Q}{\mathbb{Q}}
\newcommand{\R}{\mathbb{R}}
\newcommand{\C}{\mathbb{C}}
\newcommand{\F}{\mathbb{F}}
\newcommand{\A}{\mathbb{A}}
\newcommand{\PP}{\mathbb{P}}
\newcommand{\E}{\mathbb{E}}
\newcommand{\eps}{\varepsilon}
\newcommand{\dd}{\,\mathrm{d}}
\begin{document}
\section*{UnsolvedMath ID: 30001573}
\textbf{Source problem code:} catalogue entry
\subsection{Definitions and mathematical statement}
The supplied question does not specify the hypotheses on \(X\), the ground field, or the allowed ideals. Consequently, its precise scope remains unresolved. The following explains the usual valuation-theoretic interpretation conditionally, in a setting where function fields, birational models, log discrepancies, and the valuation formula for log canonical thresholds are defined.

A graded sequence of ideals is a family \(\mathfrak a_\bullet=(\mathfrak a_m)_{m\geq 1}\) satisfying
\[
\mathfrak a_m\mathfrak a_n\subseteq\mathfrak a_{m+n}\qquad(m,n\geq 1).
\]
For an integral \(X\), a real valuation centered on \(X\) is a valuation \(\nu:K(X)^\times\to\R\), trivial on the ground field when a ground field is specified, whose valuation ring dominates the local ring of its center on \(X\). For a nonzero finitely generated ideal \(\mathfrak b\) at that center, let \(\nu(\mathfrak b)\) be the minimum of the values of a set of local generators. For a sequence for which these values are defined, put
\[
\nu(\mathfrak a_\bullet)=\inf_{m\geq 1}\frac{\nu(\mathfrak a_m)}{m}.
\]
A valuation is quasi-monomial if, on a suitable regular birational model of \(X\), it is monomial in regular parameters cutting out the components of a simple normal crossings divisor at a stratum, with positive real weights on the selected parameters. Write \(A_X(\nu)\) for its log discrepancy over \(X\).

In a setting where the usual valuation formula applies, the weighted asymptotic log canonical threshold is expressed as
\[
\operatorname{lct}_{\mathfrak q}(\mathfrak a_\bullet)
=\inf_{\nu}\frac{A_X(\nu)+\nu(\mathfrak q)}{\nu(\mathfrak a_\bullet)},
\]
with the valuation class and conventions for zero or infinite values determined by that setting. For a valuation with \(0<\nu(\mathfrak a_\bullet)<\infty\) and finite numerator, saying that it \emph{computes} this threshold means
\[
\operatorname{lct}_{\mathfrak q}(\mathfrak a_\bullet)
=\frac{A_X(\nu)+\nu(\mathfrak q)}{\nu(\mathfrak a_\bullet)}.
\]

The source question is: for every graded sequence \(\mathfrak a_\bullet\) on \(X\) and every fixed ideal \(\mathfrak q\) in the intended, but unspecified, class, does there exist a quasi-monomial valuation \(\nu\) computing \(\operatorname{lct}_{\mathfrak q}(\mathfrak a_\bullet)\)? The missing ambient hypotheses, admissibility conditions, and boundary conventions must be recovered before this becomes a fully specified universal statement. The question asks for existence of such a valuation; it does not state that every minimizing valuation is quasi-monomial.
\subsection{Short English statement}
The question asks whether the weighted asymptotic log canonical threshold of every admissible graded sequence of ideals is attained by a quasi-monomial valuation. The supplied record leaves the geometric setting and degenerate cases unspecified. Its supplied dated assessment reports partial results, including an unweighted case, and describes a general strong statement for arbitrary \(\mathfrak q\) as open; neither that assessment nor its relationship to the displayed existence question is independently verified here.
\subsection{Sources}
\begin{itemize}
\item Local UnsolvedMath registry, permanent ID \(30001573\), source problem code \texttt{catalogue entry}, title ``Quasimonomial Valuations Computing Asymptotic Log Canonical Thresholds.'' The supplied record belongs to the Oberwolfach Reports Open Problems collection and gives proposed year 2010; no individual proposer is supplied.
\item Supplied original report citation: \emph{Komplexe Analysis (2011)}, Oberwolfach Reports, DOI 10.4171/owr/2010/39, pp.\ 2283--2333. \url{https://doi.org/10.4171/owr/2010/39}. This is also the supplied literature-source link labeled ``original OWR report.'' The report title's year differs from the year in the DOI and proposed-year metadata; the supplied citation is retained without resolution.
\item Supplied literature source: Xu, ``A minimizing valuation is quasi-monomial.'' \url{https://annals.math.princeton.edu/2020/191-3/p06}.
\item Supplied literature source: Jonsson--Mustata, ``asymptotic invariants.'' \url{https://www.numdam.org/item/10.5802/aif.2746.pdf}. The author string has corrupted diacritic markup in the record and is rendered here without accents.
\item Literature assessment dated exactly \texttt{2026-08-21}: the supplied curation labels the problem \texttt{partially\_solved} and reports quasi-monomial minimizing valuations in important cases, including the unweighted case, while describing the general strong statement for arbitrary \(\mathfrak q\) as open. This is supplied dated curation, not an independently verified literature assessment. The background expressly requests independent verification of time-sensitive status and cited claims.
\item The record labels the statement \texttt{corrected\_verified} and reports that its mathematical content was checked against the official Oberwolfach report, with notation, context, and question form normalized. This is a supplied verification claim; no source visit or independent verification is asserted here.
\end{itemize}
\end{document}
